%% ====================================================================
%% Modèle de rapport au format article, deux colonnes
%% Méthode de recherche en informatique (MRI), FASI, UPC
%% Travail pratique n° 2 : comparaison de performance entre technologies
%%
%% Compilation : pdfLaTeX (Overleaf : laisser le compilateur par défaut)
%%   pdflatex modele_article_MRI.tex   (deux fois, pour les renvois)
%% Les lignes qui commencent par % sont des commentaires : elles
%% n'apparaissent pas dans le PDF. Remplacez tout le texte entre
%% crochets [ ... ] par votre propre contenu.
%% ====================================================================
\documentclass[conference,a4paper]{IEEEtran}

\usepackage[T1]{fontenc}
\usepackage[utf8]{inputenc}
\usepackage[french]{babel}
\usepackage{mathptmx}
\usepackage{graphicx}
\usepackage{booktabs}
\usepackage{amsmath}
\usepackage{url}
\usepackage{pgfplots}
\pgfplotsset{compat=1.18}
\usepackage[hidelinks]{hyperref}

% Le modèle IEEE place la légende des tableaux en petites capitales :
% on conserve ce comportement par défaut.

\renewcommand{\IEEEkeywordsname}{Mots-clés}
\renewcommand{\abstractname}{Résumé}

\begin{document}

%% --------------------------------------------------------------------
%% TITRE : choisissez un titre précis et intéressant.
%% Exemple de forme : « Rust, Go, Node.js ou PHP : que vaut vraiment
%% la performance d'un service web sous forte charge ? »
%% --------------------------------------------------------------------
\title{[Titre de votre étude]}

%% --------------------------------------------------------------------
%% AUTEURS : un bloc par membre du groupe (nom et adresse électronique).
%% Copiez ou supprimez des blocs selon la taille du groupe.
%% \and sépare les blocs ; le modèle les répartit sur plusieurs lignes.
%% --------------------------------------------------------------------
\author{
\IEEEauthorblockN{[Prénom NOM Postnom]}
\IEEEauthorblockA{FASI, UPC\\
{[adresse@email.com]}}
\and
\IEEEauthorblockN{[Prénom NOM Postnom]}
\IEEEauthorblockA{FASI, UPC\\
{[adresse@email.com]}}
\and
\IEEEauthorblockN{[Prénom NOM Postnom]}
\IEEEauthorblockA{FASI, UPC\\
{[adresse@email.com]}}
\and
\IEEEauthorblockN{[Prénom NOM Postnom]}
\IEEEauthorblockA{FASI, UPC\\
{[adresse@email.com]}}
\and
\IEEEauthorblockN{[Prénom NOM Postnom]}
\IEEEauthorblockA{FASI, UPC\\
{[adresse@email.com]}}
\and
\IEEEauthorblockN{[Prénom NOM Postnom]}
\IEEEauthorblockA{FASI, UPC\\
{[adresse@email.com]}}
}

\maketitle

%% --------------------------------------------------------------------
%% RÉSUMÉ (150 à 250 mots) : il se lit seul. En quelques phrases :
%% le problème, ce que vous avez fait, comment, ce que vous avez trouvé,
%% ce que cela signifie. Il se rédige en dernier.
%% --------------------------------------------------------------------
\begin{abstract}
[Contexte et problème en une ou deux phrases. Objectif de l'étude.
Technologies comparées et application retenue. Démarche suivie et
métriques mesurées. Principaux résultats, chiffrés. Conclusion et portée
de ces résultats.]
\end{abstract}

\begin{IEEEkeywords}
[mot-clé 1], [mot-clé 2], [mot-clé 3], [mot-clé 4]
\end{IEEEkeywords}

%% ====================================================================
\section{Introduction}
%% ====================================================================
% Partez du général pour aller vers votre question précise.
[Contexte : pourquoi la performance d'une application compte-t-elle,
pour qui, dans quelles situations ?]

[Problème : quelles affirmations circulent sur la performance des
langages et des frameworks ? Sur quoi reposent-elles ? Citez les
travaux que vous avez trouvés, par exemple~\cite{ref1}.]

[Question de recherche : formulez-la clairement, en une phrase.]

[Hypothèses : ce que vous vous attendez à observer, et pourquoi.]

[Contribution et plan : ce que cet article apporte, puis l'organisation
des sections suivantes.]

%% ====================================================================
\section{Méthodologie}
%% ====================================================================
\subsection{Application de référence}
[Décrivez l'application et son cahier des charges : fonctionnalités,
données manipulées, comportement attendu.]

\subsection{Technologies comparées}
[Présentez au moins quatre technologies et justifiez leur choix.
Indiquez les versions utilisées.]

\begin{table}[htbp]
\caption{Technologies comparées}
\label{tab:techno}
\centering
\begin{tabular}{@{}lll@{}}
\toprule
\textbf{Technologie} & \textbf{Langage} & \textbf{Version} \\
\midrule
{[Technologie 1]} & [Langage] & [x.y.z] \\
{[Technologie 2]} & [Langage] & [x.y.z] \\
{[Technologie 3]} & [Langage] & [x.y.z] \\
{[Technologie 4]} & [Langage] & [x.y.z] \\
\bottomrule
\end{tabular}
\end{table}

\subsection{Métriques}
[Quelles grandeurs mesurez-vous pour juger qu'une version est plus
performante qu'une autre ? Pour chacune : définition, unité, outil de
mesure, raison de son choix.]

\subsection{Outils et environnement de test}
[Outil de test de charge, outils de mesure, machine utilisée.]

\subsection{Déroulement des tests}
[Décrivez précisément ce que vous avez fait, dans l'ordre, de façon à
ce qu'un autre groupe puisse refaire les mêmes tests.]

%% ====================================================================
\section{Développement}
%% ====================================================================
\subsection{Réalisation des versions}
[Comment chaque version a-t-elle été écrite ? Structure du code,
bibliothèques utilisées, points particuliers.]

\subsection{Équivalence fonctionnelle}
[Comment avez-vous vérifié que toutes les versions font exactement la
même chose et renvoient les mêmes résultats ?]

\subsection{Difficultés rencontrées}
[Problèmes techniques et solutions adoptées.]

%% ====================================================================
\section{Résultats et discussion}
%% ====================================================================
\subsection{Données recueillies}
[Présentez les données : volume, format, tableau de synthèse.]

\begin{table}[htbp]
\caption{Synthèse des mesures (à compléter)}
\label{tab:resultats}
\centering
\begin{tabular}{@{}lccc@{}}
\toprule
\textbf{Technologie} & \textbf{[Métrique 1]} & \textbf{[Métrique 2]} & \textbf{[Métrique 3]} \\
 & [unité] & [unité] & [unité] \\
\midrule
{[Technologie 1]} & ... & ... & ... \\
{[Technologie 2]} & ... & ... & ... \\
{[Technologie 3]} & ... & ... & ... \\
{[Technologie 4]} & ... & ... & ... \\
\bottomrule
\end{tabular}
\end{table}

\subsection{Analyse}
[Décrivez ce que montrent les données et les graphiques. Chaque figure
est citée dans le texte, par exemple : la figure~\ref{fig:exemple}
montre que ...]

% Exemple de graphique en barres. Les valeurs sont FICTIVES :
% remplacez-les par vos propres mesures (ou importez vos données
% depuis un fichier CSV avec \addplot table).
\begin{figure}[htbp]
\centering
\begin{tikzpicture}
\begin{axis}[
  width=\columnwidth, height=5cm,
  ybar, bar width=12pt, ymin=0,
  symbolic x coords={Techno 1,Techno 2,Techno 3,Techno 4},
  xtick=data,
  ylabel={[Métrique] ([unité])},
  nodes near coords,
  nodes near coords style={font=\scriptsize},
  tick label style={font=\footnotesize},
  label style={font=\footnotesize},
  enlarge x limits=0.18]
\addplot[fill=gray!40, draw=black] coordinates
  {(Techno 1,12) (Techno 2,18) (Techno 3,25) (Techno 4,40)};
\end{axis}
\end{tikzpicture}
\caption{[Titre du graphique] (valeurs fictives, à remplacer).}
\label{fig:exemple}
\end{figure}

% Exemple de courbes : évolution d'une métrique selon la charge.
\begin{figure}[htbp]
\centering
\begin{tikzpicture}
\begin{axis}[
  width=\columnwidth, height=5cm,
  xlabel={[Niveau de charge]},
  ylabel={[Métrique] ([unité])},
  tick label style={font=\footnotesize},
  label style={font=\footnotesize},
  legend style={font=\scriptsize, at={(0.03,0.97)}, anchor=north west},
  ymin=0]
\addplot[black, thick, mark=*] coordinates {(10,2) (50,3) (100,5) (200,8)};
\addplot[black, thick, dashed, mark=square*] coordinates {(10,3) (50,5) (100,8) (200,14)};
\addplot[black, thick, dotted, mark=triangle*] coordinates {(10,5) (50,9) (100,15) (200,30)};
\addplot[gray, thick, dashdotted, mark=diamond*] coordinates {(10,9) (50,16) (100,28) (200,60)};
\legend{Techno 1, Techno 2, Techno 3, Techno 4}
\end{axis}
\end{tikzpicture}
\caption{[Titre du graphique] (valeurs fictives, à remplacer).}
\label{fig:courbes}
\end{figure}

\subsection{Interprétation et discussion}
[Comment expliquez-vous les écarts observés ? Les résultats
confirment-ils vos hypothèses ? Que ne peut-on pas conclure à partir de
ces données ? Quelles sont les limites de votre étude ?]

%% ====================================================================
\section{Conclusion}
%% ====================================================================
[Rappel de la question, réponse apportée, limites, suites possibles.]

%% ====================================================================
%% RÉFÉRENCES : citez uniquement des sources réelles, que vous avez
%% consultées. Dans le texte : \cite{ref1}.
%% ====================================================================
\begin{thebibliography}{00}
\bibitem{ref1} [Auteur(s)], « [Titre de l'article] », \textit{[Revue ou
conférence]}, vol.~[x], n\textsuperscript{o}~[x], p.~[x-y], [année].
\bibitem{ref2} [Auteur(s)], \textit{[Titre de l'ouvrage]}. [Ville] :
[Éditeur], [année].
\bibitem{ref3} [Auteur ou organisation], « [Titre de la page] », [année].
[En ligne]. Disponible : \url{https://...}. [Consulté le jj mois aaaa].
\end{thebibliography}

\end{document}
